\section{The allocation problem before the optimizer}
\label{sec:introduction}

An electrical drive does not ask for three currents; it asks for torque while
respecting an inverter and a thermal system. In an \gls{eesm}, infinitely many
triples \((\gls{sym:id},\gls{sym:iq},\gls{sym:iexc})\) can produce the same
torque. The current-reference problem is therefore an allocation problem: how
should the demanded electromagnetic effect be shared between stator saliency,
stator quadrature current, and rotor excitation? The corresponding \gls{psm}
problem is the two-dimensional fixed-excitation slice of that question.

The usual route writes a Lagrangian immediately. That route is correct but it
conceals the reason small solutions exist. Before differentiating anything,
consider the three surfaces in current space. Equal copper loss is a closed
ellipsoid. Equal nonzero torque is a hyperbolic cylinder. Equal stator-voltage
magnitude is an oblique elliptic cylinder because three current inputs feed
only two stator-voltage outputs. Minimum loss means expanding the loss
ellipsoid until it first touches a requested torque surface without leaving
the voltage cylinder. Maximum torque means moving through the same nested
surfaces until the first physical boundary is reached. The two tasks are thus
different questions about one operating frontier, not two unrelated
optimizers.

Two precise problems will be connected. For requested torque
\gls{sym:mref}, define
\begin{equation}
 \gls{sym:preq}(\gls{sym:mref})
 =\min_{\gls{sym:i}}\left\{\gls{sym:pcu}(\gls{sym:i})\;\middle|\;
 \gls{sym:torque}(\gls{sym:i})=\gls{sym:mref},\ 
 \norm{\gls{sym:u}(\gls{sym:i})}\leq\gls{sym:umax}\right\}.
 \label{eq:preq-definition}
\end{equation}
The requested point is physically feasible only if
\begin{equation}
 \boxed{\gls{sym:preq}(\gls{sym:mref})\leq\gls{sym:pcumax}.}
 \label{eq:feasibility-test}
\end{equation}
The loss ceiling is deliberately not included as a redundant constraint in
the minimization: it is the acceptance test on the minimum. Removing it from
the controller would remove the thermal feasibility statement, not simplify
the physics. The companion problem is
\begin{equation}
 \gls{sym:mmax}=\max_{\gls{sym:i}}\left\{\gls{sym:torque}(\gls{sym:i})
 \;\middle|\;\gls{sym:pcu}\leq\gls{sym:pcumax},\ 
 \norm{\gls{sym:u}}\leq\gls{sym:umax}\right\}.
 \label{eq:mmax-definition}
\end{equation}

The same value functions will later be read in reverse. Instead of fixing a
machine and selecting a current, fix a required torque-speed envelope and ask
which electromagnetic parameter vectors admit any feasible current. This
creates a machine-design region whose boundary is the level set
\(M_{\max}(\boldsymbol\theta)=M_{\mathrm{req}}\). The forward and inverse
questions therefore share one solver and one geometry; only the variable being
projected out changes.

A third question changes the mathematical object from a point to a curve.
Given two feasible currents, what voltage-constrained trajectory reaches the
target fastest, with least transient energy, or with a declared compromise?
The optimal steady endpoint and the optimal route to it are distinct. The
dynamic part therefore restores Faraday's law, reciprocal rotor-field
dynamics, converter geometry, and electrical drift before defining distance.

A fourth question precedes all three mathematically even though it is
developed after them didactically: in which coordinates should sensing,
reference generation, optimization, and control be performed? The paper
therefore distinguishes spatial electromagnetic frames, current-space task
charts, and virtual input coordinates. Rotor \(dq\), stator-flux \(ft\),
active flux, torque/performance coordinates, and local metric frames are
judged by what they simplify and by the observer, inverse, conditioning, and
real-time costs they move elsewhere.

The paper's organizing observation is homogeneity. In the linear EESM model,
loss, torque, and squared voltage all grow with the square of current scale.
A current direction therefore chooses three efficiencies; one scalar radius
then chooses the operating level. This separation leads to three mutually
illuminating descriptions:
\begin{itemize}
 \item loss whitening and torque-aligned coordinates expose the physical
       factors that multiply to create torque;
 \item hyperbolic, projective, and conformal coordinates explain the geometry
       of a constant product;
 \item the performance plane maps every current direction directly to voltage
       cost and torque benefit per unit copper loss.
\end{itemize}
The last description supplies the preferred controller core: a one-dimensional
support-direction search involving symmetric \(3\times3\) eigenproblems. The
other descriptions explain why that reduction works and provide independent
checks. KKT conditions and algebraic resultants are retained as reference
routes, following the established EESM setpoint literature
\cite{reinhard2022} and the quadrics-and-quartics treatment of anisotropic
synchronous machines \cite{eldeeb2018,eldeeb2016}.

The development is intentionally self-contained. The PSM is first used as a
two-dimensional microscope: its exact resistive voltage ellipse can be drawn,
translated, and rotated. The three-dimensional EESM is then easier to see as
a homogeneous parent problem. The relation to established active flux is
stated carefully: active flux motivates the torque-producing combination
\cite{boldea2008}; the subsequent loss-orthogonal completion and projective
geometry are coordinate developments for the allocation problem, not a claim
to rename that literature.

The initial operating-point model is linear, stationary, and peak-valued.
Saturation, cross-saturation, iron loss, inverter nonlinearity, and current
dynamics are not hidden inside that derivation. Section~\ref{sec:nonlinear}
adds nonlinear magnetics through a local affine model, while Part~III derives
the full current dynamics and their state-dependent local geometry. No
measured machine data were found
in the workspace; numerical examples are therefore explicitly synthetic and
serve as regression tests, never as experimental validation.

\paragraph{Reading strategy.}
Sections~\ref{sec:model}--\ref{sec:psm-ellipse} identify the surfaces.
Sections~\ref{sec:frontier}--\ref{sec:performance} discover useful
coordinates. Sections~\ref{sec:algebraic}--\ref{sec:controller} translate the
geometry into verification and implementation. The limiting cases and final
comparison are separated so that a compact formula is never mistaken for a
universally valid one. Part II then turns operating value functions into
loss-limited design maps, voltage-aware level surfaces, and intersections for
multiple operating points. Part III develops dynamic operating-point
transitions from reciprocal flux dynamics through velocity sets, Finsler
navigation, reachability, and a reproducible fast--slow EESM experiment.
Part IV develops the representation problem from \(dq\) and DFVC through
active flux, local task charts, feedback linearization, conditioning maps, and
embedded operation counts. Part V distinguishes electromagnetic parameter
space from the non-unique preimage of physical machine geometries and closes
the loop to machine-control co-design.
